\documentclass{article}

% Keep the official ICLR files in ../template/iclr2026/ while making this
% manuscript directly compilable from the main/ directory.
\makeatletter
\def\input@path{{../template/iclr2026/}}
\makeatother
\usepackage{iclr2026_conference,times}
\input{math_commands.tex}

\usepackage{fontspec}
\usepackage{xeCJK}
\setmainfont{TeX Gyre Termes}
\setCJKmainfont{AR PL UMing CN}
\setCJKsansfont{Droid Sans Fallback}
\setCJKmonofont{Droid Sans Fallback}
\XeTeXlinebreaklocale "zh"
\XeTeXlinebreakskip=0pt plus 1pt

\usepackage{amsmath,amssymb,mathtools}
\usepackage{booktabs,tabularx,threeparttable,multirow,array}
\usepackage{graphicx,xcolor,enumitem,float}
\usepackage{tikz}
\usetikzlibrary{arrows.meta,positioning,fit,calc,shapes.geometric,backgrounds}
\usepackage{pgfplots}
\pgfplotsset{compat=1.18}
\usepackage{hyperref}
\usepackage{url}

\hypersetup{colorlinks=true,linkcolor=blue,citecolor=blue,urlcolor=blue}
\setlist{nosep,leftmargin=*}
\newcolumntype{Y}{>{\centering\arraybackslash}X}
\newcommand{\scc}{\textsc{SC-CL}}
\newcommand{\bone}{\textsc{B1}}

\title{开放世界知识图谱关系候选排序的可审计结构反事实对比学习}
\author{}

% \iclrfinalcopy % Keep commented for the anonymous review template.
\begin{document}

\maketitle

\begin{abstract}
生成式知识图谱关系预测需要能够揭示关系混淆、且来源可追溯的训练候选。开放世界条件下，候选的排序难度与事实真值需要分别处理。本文提出结构反事实对比学习（Structure-Counterfactual Contrastive Learning, \scc），融合训练图的关系结构画像与冻结 GRIP 基线的分数违例度，在过滤实体对训练已知关系后构建固定候选库，并以候选级 InfoNCE 联合优化关系排序与原始生成目标。在 NELL23K 的单种子评估中，\scc{} 在 4,943 条测试问题上达到 91.52\% 的候选重排 H@1 和 95.27\% 的 MRR；在同候选列表的前 200 条配对问题上，MRR 和 H@1 分别比基线高 1.60 和 2.50 个百分点。\scc{} 将结构与模型混淆信号纳入同一候选构造流程，为开放世界关系候选学习提供可追溯的训练接口。
\end{abstract}

\section{引言}

给定实体对及其图上下文，关系预测需要识别相应的关系标签。GRIP 将图结构融入语言模型参数，通过生成式问答完成这一任务 \citep{feng2025grip}。候选级对比学习让目标关系与替代关系直接竞争，而候选的选取决定模型接收到什么关系辨别信号。因此，候选构造是连接图结构与关系排序目标的关键环节。

困难候选的构造需要同时处理结构相关性、模型混淆与事实不确定性。训练图的关系统计描述不同关系的结构接近程度，模型分数则揭示已有的关系混淆。知识图谱的不完整性进一步要求区分候选的学习用途与事实真值：未标注关系在开放世界中保留未判定状态。本文研究如何将结构与模型混淆信号共同用于关系候选构造，并将训练已知关系过滤和来源追溯纳入同一流程。

为此，我们提出结构反事实对比学习（\scc）。\scc{} 以关系结构画像刻画候选在训练图中的统计接近程度，以冻结的 generation-only 基线 \bone{} 的分数违例度刻画模型混淆，两种信号共同决定带均匀下限的采样分布。训练已知关系过滤与固定候选库明确候选来源，使每个候选能够关联到其构造记录。\scc{} 接入 GRIP Stage~2，联合优化候选级 InfoNCE 与原始生成目标，沿用 Stage~1 和原始提示。本文将结构反事实定义为固定实体对上的关系标签替换，图~\ref{fig:mechanism} 展示其构造与学习流程。

在 NELL23K 的 4,943 条测试问题上，\scc{} 的候选重排 H@1 为 91.52\%，MRR 为 95.27\%。在相同候选列表的前 200 条配对问题上，其 MRR 比 \bone{} 高 1.60 个百分点，H@1 高 2.50 个百分点。实验节分别展示完整测试集的排序水平和配对子集的逐题预测变化，并报告直接解码结果、候选来源检查及分阶段成本。评估设置与比较范围见第~\ref{sec:experiments} 节。

本文面向固定实体对上的关系候选排序，贡献包括三个方面：定义结构画像与冻结模型违例度联合驱动的训练候选分布；构建包含训练已知关系过滤、固定候选库和逐候选来源记录的学习流程；在 NELL23K 展示全量候选重排结果与同候选前 200 题的配对分析。

\section{相关工作}

\paragraph{关系预测与生成式图推理。}
TransE、ComplEx、RotatE 和 ConvE 通过不同的三元组打分函数建模实体与关系 \citep{bordes2013transe,trouillon2016complex,sun2019rotate,dettmers2018conve}。KG-BERT 将三元组判断转为文本建模 \citep{yao2019kgbert}，SimKGC 在预训练语言模型上使用多种实体负例进行对比学习 \citep{wang2022simkgc}。GRIP 以图上下文和生成式问答学习参数化图推理 \citep{feng2025grip}。本文关注固定实体对上的关系标签竞争，在 GRIP 的生成式训练框架中研究候选来源与候选级对比目标的结合。

\paragraph{困难候选与假阴性。}
结构感知负采样 SANS 和对抗式采样 KBGAN 分别从图结构与模型反馈寻找困难负例 \citep{ahrabian2020sans,cai2018kbgan}。去偏对比学习与 ProGCL 讨论了对比学习中候选难度和假阴性风险的关系 \citep{chuang2020dcl,xia2022progcl}。\scc{} 将训练图关系画像与冻结模型的分数违例度共同用于关系候选构造，并以训练已知关系过滤、均匀采样下限和固定候选库明确候选的构造规则。本文采用的开放世界真值约定见问题定义节。

\paragraph{关系候选评估。}
传统 link prediction 常在实体替换候选上报告 MRR/Hits，CoDEx 还提供经过人工核验的实体替换困难负例 \citep{safavi2020codex}。实体替换与关系替换对应不同的候选空间。本文在固定实体对上评价关系标签排序，并分别使用直接解码与逐候选评分衡量两种答案选择方式；实体替换负例保留其三元组判断用途，关系候选按各自的已知事实与未判定状态记录。

\section{问题定义}

\subsection{训练图与候选关系}

令训练知识图谱为 $\mathcal{G}_{\mathrm{tr}}=(\mathcal{V},\mathcal{E}_{\mathrm{tr}})$，其中
$\mathcal{E}_{\mathrm{tr}}\subseteq \mathcal{V}\times\mathcal{R}_{\mathrm{tr}}\times\mathcal{V}$，
$\mathcal{R}_{\mathrm{tr}}$ 是训练关系词表。一个关系问答样本表示为
$q=(h,t,r^+)$，其中 $r^+$ 是目标关系，$h,t$ 是有序实体对。定义同一实体对在训练图中的已知关系集合为
\begin{equation}
    T_{\mathrm{tr}}(h,t)=\{r:(h,r,t)\in\mathcal{E}_{\mathrm{tr}}\}.
\end{equation}
候选池为
\begin{equation}
    P_q=\mathcal{R}_{\mathrm{tr}}\setminus\bigl(\{r^+\}\cup T_{\mathrm{tr}}(h,t)\bigr).
    \label{eq:pool}
\end{equation}
式~\ref{eq:pool} 从训练关系词表中排除目标关系和该实体对的已知训练关系。其余关系作为真值未判定的训练替代项，遵循知识图谱的开放世界假设 \citep{ji2022survey,hogan2021kg}。本文的结构反事实操作固定实体对 $(h,t)$，将目标关系替换为 $r\in P_q$；“负候选”表示其在对比目标中的角色，事实真值保持未判定。

\subsection{答案选择与候选学习目标}

本文区分两种答案选择规则。直接解码在列出十个关系选项的题目上输出关系文本，以 Exact Match（EM）评价原始解码路径下的预测。候选重排给定同一题目的候选集合，用关系续接分数（continuation score）$z_\theta(q,r)$ 逐项排序，以 Hits@1（H@1）、H@3 和 MRR 评价候选间的关系辨别。重排 H@1 等于按最高分选项计算的 EM。两种规则分别报告，模型间的表现差异采用相同答案选择规则比较。

给定冻结 bank $B(q)=\{r_1^-,\ldots,r_K^-\}$，候选损失写为
\begin{equation}
    \mathcal{L}_{\mathrm{cand}}(q)=-\log
    \frac{\exp(z_\theta(q,r^+)/\tau_c)}
    {\exp(z_\theta(q,r^+)/\tau_c)+\sum_{k=1}^K\exp(z_\theta(q,r_k^-)/\tau_c)}.
    \label{eq:infonce}
\end{equation}
其中 $\tau_c$ 是候选温度。训练阶段联合优化原始生成损失 $\mathcal{L}_{\mathrm{gen}}$ 与候选损失：
\begin{equation}
    \mathcal{L}=\mathcal{L}_{\mathrm{gen}}+\lambda_{\mathrm{cand}}\mathcal{L}_{\mathrm{cand}}.
    \label{eq:total}
\end{equation}
\scc{} 沿用 Stage~1 和生成目标，将新增训练项明确为候选级 InfoNCE。

\section{结构反事实对比学习}

\scc{} 将候选构造与 Stage~2 优化分为两个阶段（图~\ref{fig:mechanism}）。离线阶段融合训练图关系画像与冻结 B1 在训练问题上的违例度，在过滤后的关系池中采样并固定每题九个候选。学习阶段从 Stage~1 adapter 初始化 GRIP，在固定候选库上联合优化候选损失与原始生成损失。

\begin{figure}[H]
    \centering
    \resizebox{\linewidth}{!}{\begin{tikzpicture}[
    font=\small,
    box/.style={draw=black!35, rounded corners=3pt, align=center,
                minimum height=10mm, text width=29mm, fill=white},
    signal/.style={box, draw=teal!65!black, fill=teal!6},
    train/.style={box, draw=blue!60!black, fill=blue!5},
    flow/.style={-{Latex[length=2mm]}, line width=0.8pt, draw=black!65},
    offline/.style={flow, dashed},
    panel/.style={rounded corners=5pt, line width=0.8pt}
]
\begin{scope}[on background layer]
    \draw[panel, draw=teal!50, fill=teal!2] (-1.9,1.0) rectangle (5.9,-5.5);
    \draw[panel, draw=blue!40, fill=blue!2] (6.25,1.0) rectangle (13.75,-5.5);
\end{scope}
\node[anchor=west, text=teal!65!black, font=\bfseries] at (-1.65,0.65)
    {A\quad 离线候选构造};
\node[anchor=west, text=blue!65!black, font=\bfseries] at (6.5,0.65)
    {B\quad Stage 2 联合学习};

\node[signal] (graph) at (0, -0.35) {训练图与关系词表\\$\mathcal G_{\rm tr},\mathcal R_{\rm tr}$};
\node[signal] (profile) at (0,-1.85) {五维关系结构画像\\$s_{\rm str}(q,r)$};
\node[signal] (b1) at (4,-0.35) {冻结 B1\\训练问题离线评分};
\node[signal] (violation) at (4,-1.85) {分数违例度\\$V_{\rm B1}(q,r)$};
\node[signal, text width=61mm] (pool) at (2,-3.25)
    {train-known 过滤与双信号混合采样\\
     $u=\alpha s_{\rm str}+\beta V_{\rm B1}$，均匀下限 $\rho$};
\node[signal, text width=61mm] (bank) at (2,-4.75)
    {固定候选库：每题 $K=9$，无放回采样\\来源、信号、过滤记录与文件哈希};
\draw[flow] (graph) -- (profile);
\draw[offline] (b1) -- (violation);
\draw[flow] (profile.south) -- (pool.north west);
\draw[flow] (violation.south) -- (pool.north east);
\draw[flow] (pool) -- (bank);

\node[train] (s1) at (8.1,-0.35) {Stage 1 adapter\\参数初始化};
\node[train] (qa) at (11.9,-0.35) {原始训练问题\\图上下文与目标关系};
\node[train, text width=59mm] (model) at (10,-1.85)
    {GRIP：冻结基座模型 + 可训练 LoRA};
\node[train] (gen) at (8.1,-3.25) {生成损失\\$\mathcal L_{\rm gen}$};
\node[train] (cand) at (11.9,-3.25) {候选评分与 InfoNCE\\$\mathcal L_{\rm cand}$};
\node[train, text width=59mm] (total) at (10,-4.75)
    {$\mathcal L=\mathcal L_{\rm gen}+\lambda_{\rm cand}\mathcal L_{\rm cand}$\\联合更新 LoRA 参数};
\draw[flow] (s1.south) -- (model.north west);
\draw[flow] (qa.south) -- (model.north east);
\draw[flow] (model.south west) -- (gen.north);
\draw[flow] (model.south east) -- (cand.north);
\draw[flow] (gen.south) -- (total.north west);
\draw[flow] (cand.south) -- (total.north east);
\draw[offline] (bank.east) -- (5.65,-4.75) -- (5.65,-2.55)
    -- node[above, font=\scriptsize, fill=white, inner sep=2pt] {训练期固定候选}
    (11.9,-2.55) -- (cand.north);
\end{tikzpicture}
}
    \caption{\scc{} 的候选构造与联合学习。A：训练图结构画像与冻结 \bone{} 违例度共同指导采样，train-known 过滤与无放回采样形成每题 $K=9$ 的固定候选库。B：Stage~2 从 Stage~1 adapter 初始化，联合优化生成损失与加权候选 InfoNCE；冻结 \bone{} 为离线候选构造提供评分。}
    \label{fig:mechanism}
\end{figure}

\subsection{训练图关系画像}

对每一个 $r\in\mathcal{R}_{\mathrm{tr}}$，我们从 $\mathcal{E}_{\mathrm{tr}}$ 计算五维关系画像
\begin{equation}
    \phi(r)=\bigl(\log(1+|E_r|),\log(1+|H_r|),\log(1+|T_r|),
    \overline{d}_{H,r},\overline{d}_{T,r}\bigr),
    \label{eq:profile}
\end{equation}
其中 $|E_r|$ 为关系边数，$|H_r|$ 和 $|T_r|$ 分别为关系覆盖的头实体数与尾实体数，
$\overline{d}_{H,r}$ 是关系边上头实体的平均出度，$\overline{d}_{T,r}$ 是尾实体的平均入度。边数与实体覆盖计数使用 $\log(1+x)$ 压缩数值尺度；度统计保留图的局部连接性。

对目标关系 $r^+$ 和候选 $r$，定义归一化结构相似度
\begin{equation}
    s_{\mathrm{str}}(q,r)=1-\frac{1}{5}\sum_{j=1}^{5}
    \frac{|\phi_j(r^+)-\phi_j(r)|}
    {|\phi_j(r^+)|+|\phi_j(r)|+\epsilon},
    \label{eq:similarity}
\end{equation}
其中 $\epsilon$ 保证分母非零，分数裁剪到 $[0,1]$。该分数衡量五维图统计画像的接近程度；实体对的已知关系由式~\ref{eq:pool} 单独过滤。

\subsection{冻结 B1 的边界违例度}

\bone{} 是原始 GRIP 的 generation-only 基线。我们在训练问题上离线计算全关系词表的 continuation score。若离线记录已有 pairwise confusion，则直接使用裁剪到 $[0,1]$ 的 confusion；否则对候选和正关系的分数差使用平滑违例度
\begin{equation}
    V_{\bone}(q,r)=\sigma\left(\frac{z_{\bone}(q,r)-z_{\bone}(q,r^+)}{\tau_v}\right).
    \label{eq:violation}
\end{equation}
候选分数越接近或超过正关系，$V_{\bone}$ 越大。B1 在训练问题上的离线分数于候选库构造时固定，提供与 Stage~2 参数更新分离的候选评分依据。

\subsection{结构反事实候选分布}

将结构相似度和分数违例度合成为
\begin{equation}
    u(q,r)=\alpha s_{\mathrm{str}}(q,r)+\beta V_{\bone}(q,r),
    \qquad r\in P_q.
    \label{eq:utility}
\end{equation}
候选分布使用均匀下限与 utility softmax 的混合：
\begin{equation}
    p(r\mid q)=\frac{\rho}{|P_q|}+(1-\rho)
    \frac{\exp(u(q,r)/\tau_s)}{\sum_{r'\in P_q}\exp(u(q,r')/\tau_s)}.
    \label{eq:sampling}
\end{equation}
我们从 $p(r\mid q)$ 无放回抽取 $K=9$ 个负候选。主实验参数为 $\rho=1/3$、$\alpha=1$、$\beta=1$、$\tau_s=0.25$。当 $|P_q|<K$ 时，记录该样本及候选不足原因，保持每个有效样本的候选互异性；实验报告候选构造失败数。

结构反事实操作在固定实体对上替换关系标签，其候选池与开放世界真值约定见问题定义节。

\subsection{固定候选库与来源追溯}

\scc{} 为每个训练问题记录问题标识、目标关系、九个候选及其结构相似度、违例度、utility 和 B1 score gap，并保存采样参数、数据划分与已知关系过滤记录。候选库遵循四项来源约束：

\begin{enumerate}
    \item 问题来源和已知关系过滤均使用训练划分；
    \item 每个样本恰有一个正关系和 $K=9$ 个互异候选，候选均属于 $\mathcal{R}_{\mathrm{tr}}$；
    \item 候选与 $T_{\mathrm{tr}}(h,t)$ 不相交，问题标识与训练任务一一对齐；
    \item 候选库哈希、脚本版本、随机种子和采样参数写入来源清单，构造后的候选库在训练中保持固定。
\end{enumerate}

训练时，候选 continuation score 代入式~\ref{eq:infonce}，生成损失沿用原始 GRIP 路径计算。两项损失按式~\ref{eq:total} 共同优化 Stage~2 的 LoRA 参数。

\subsection{计算复杂度与离线构造}

关系画像通过一次线性扫描训练边得到，复杂度为 $O(|\mathcal{E}_{\mathrm{tr}}|)$。给定离线模型分数后，每个问题在关系词表上计算 utility，复杂度为 $O(|\mathcal{R}_{\mathrm{tr}}|)$；采样后的候选库大小为 $O(K|Q_{\mathrm{tr}}|)$，其中 $Q_{\mathrm{tr}}$ 为训练问题集合。计算过程包括训练图统计、冻结 B1 离线评分和 Stage~2 联合训练；分阶段资源记录见第~\ref{sec:cost} 节。

\section{实验与结果}
\label{sec:experiments}

\subsection{实验设置与评价指标}

主实验使用 NELL23K、Qwen2.5-7B-Instruct 和训练种子 2026，在 4,943 条十选一关系问题上评价 \scc{}。候选重排通过 teacher-forced continuation score 逐一评分，以 H@1、MRR 和 H@3 衡量给定候选集合中的关系排序；直接解码 EM 衡量原始输出路径下的预测。两种评估使用相同题目与选项，分别报告结果。

\bone{} 为原始 GRIP 的 generation-only 基线。B1 与 \scc{} 的同候选重排比较覆盖测试文件的前 200 题；两者沿用各自现有训练配置，比较条件集中见第~\ref{sec:comparison-scope} 节。完整测试集报告 \scc{} 的排序表现，以及两个模型的直接解码表现。训练设置见附录表~\ref{tab:hyper}，WN18RR 的直接解码补充结果见附录~\ref{sec:wn-exploratory}。

\subsection{NELL23K 候选排序结果}

\scc{} 在完整测试集上达到 95.27\% MRR、91.52\% H@1 和 99.07\% H@3。在同候选前 200 题上，\scc{} 的 MRR 为 94.96\%，比 B1 高 1.60 个百分点；H@1 为 91.00\%，比 B1 高 2.50 个百分点。表~\ref{tab:nell-results} 给出全量结果与配对结果。

\begin{table}[htbp]
\caption{NELL23K 候选重排结果（\%，训练种子 2026）。配对区块按测试文件顺序取前 200 题，两个模型使用相同问题与候选列表；全量区块报告 \scc{} 的 4,943 题结果，全量 B1 重排尚未完成。$\Delta$ 为配对子集上的百分点差。}
\label{tab:nell-results}
\centering\small
\begin{tabular}{llrrrr}
\toprule
范围 & 模型 & $n$ & H@1 & MRR & H@3 \\
\midrule
配对前 200 题 & \bone & 200 & 88.50 & 93.36 & 97.50 \\
配对前 200 题 & \scc-9 & 200 & \textbf{91.00} & \textbf{94.96} & \textbf{99.50} \\
配对差异 & $\Delta$(\scc-9, \bone) & -- & +2.50 & +1.60 & +2.00 \\
\midrule
全量测试 & \scc-9 & 4,943 & 91.52 & 95.27 & 99.07 \\
\bottomrule
\end{tabular}
\end{table}

\subsection{配对预测的逐题变化}

在同候选前 200 题上，9 题由 B1 错误转为 \scc{} 正确，4 题发生反向变化，净增加 5 题正确预测；其余为 173 题共同正确和 14 题共同错误。逐题变化与 H@1 的 2.50 个百分点差异一致，并展示两模型预测差异的分布。

\subsection{直接解码与候选评分}

在相同的 4,943 条测试问题上，\scc{} 与 B1 的直接解码 EM 分别为 87.03\% 和 87.46\%，相差 $-0.43$ 个百分点（表~\ref{tab:direct-decode}）。\scc{} 通过逐候选评分选择答案时，H@1 为 91.52\%。直接解码考察关系文本输出，逐候选评分考察给定选项间的排序，本文围绕后一目标组织主要分析。

\begin{table}[htbp]
\caption{NELL23K 直接解码结果（\%，4,943 题，训练种子 2026）。模型在列出相同十个选项的题目上生成关系文本，按 Exact Match 计分。}
\label{tab:direct-decode}
\centering\small
\begin{tabular}{lr}
\toprule
模型 & 直接解码 EM \\
\midrule
\bone & 87.46 \\
\scc-9 & 87.03 \\
\bottomrule
\end{tabular}
\end{table}

\subsection{候选构造与来源追溯}

NELL23K 的固定候选库覆盖 3,253 条训练问题，每题包含九个互异关系，候选构造失败数为 0。来源检查记录了训练划分、任务问题对齐与候选库标识；采样清单保存结构画像、B1 违例度及采样参数。\scc{} Stage~2 完成 10 个 epoch、240 个优化步，归档了 adapter、直接解码预测和候选重排输出。这些记录将候选构造与模型训练关联到同一运行，便于追溯每个候选的来源。

\subsection{离线构造与运行成本}
\label{sec:cost}

训练运行记录的硬件为 NVIDIA GeForce RTX 3090（24 GiB）。表~\ref{tab:cost} 按阶段报告本次运行的资源用量：离线 B1 评分覆盖 644,094 个训练关系候选，Stage~2 处理 12,148,720 个输入 token，完整测试重排覆盖 4,943 条问题。固定候选库在训练期间复用离线采样结果。

\begin{table}[htbp]
\caption{NELL23K 单次运行的分阶段成本。显存数值对应离线评分阶段的记录。}
\label{tab:cost}
\centering\small
\begin{tabular}{lrr}
\toprule
阶段 & 耗时（秒） & 峰值显存（GiB） \\
\midrule
冻结 B1 离线评分 & 19,027 & 15.21 \\
Stage~2 联合训练 & 80,870 & -- \\
全测试候选重排 & 3,567 & -- \\
\bottomrule
\end{tabular}
\end{table}

\section{讨论}
\label{sec:comparison-scope}

\paragraph{候选构造的双信号设计。}
关系结构画像描述候选在训练图中的统计接近程度，冻结 B1 分数描述模型已有的关系混淆。\scc{} 将二者作为候选采样依据，并保存逐候选的构造记录，使进入训练的结构与混淆模式可供检查。配对子集中的排序差异与逐题变化提供了进一步分析候选学习的切入点。

\paragraph{训练目标与答案选择。}
候选 InfoNCE 显式训练目标关系与替代关系之间的分数竞争。逐候选评分直接评价这一排序目标，直接解码同时涉及关系文本输出。本文分别报告两种答案选择方式下的表现，以区分给定候选的排序能力与原始解码路径下的预测表现。

\paragraph{固定候选库的适用条件。}
固定候选库将离线构造与 Stage~2 优化分离，使同一运行中的候选来源和采样信号保持固定。其适用条件是一次冻结 B1 全关系评分的预处理成本，以及训练期间固定使用的候选分布。方法面向固定实体对上的关系标签替换，利用训练图统计和关系词表构造排序候选。

\paragraph{比较设计与验证范围。}
主实验对应 NELL23K 的训练种子 2026。B1 与 \scc{} 的同候选重排比较覆盖测试文件前 200 题，采用各自现有训练配置，训练预算存在差异；全量 4,943 题重排结果对应 \scc{}。全量 B1 重排、同预算随机及频率匹配对照、多种子重复与持出正例重叠审计尚未完成。文中的配对差异描述该子集上的模型表现，组合采样的独立作用与跨数据集适用性属于后续验证问题。候选审计按训练已知关系冲突、持出正例重叠和开放世界未判定分别记录，具体规则见附录~\ref{sec:audit-protocol}。

\section{结论}

本文提出 \scc{}，将训练图关系结构画像与冻结模型违例度结合，通过训练已知关系过滤和固定候选库，为 GRIP 的候选级对比学习提供可追溯的关系候选。在 NELL23K 的 4,943 条测试问题上，\scc{} 的候选重排达到 91.52\% H@1 和 95.27\% MRR；同候选前 200 题上的 MRR 比 B1 高 1.60 个百分点。该设计将候选的结构依据、模型混淆信号与来源追溯纳入同一训练流程，为开放世界关系候选学习提供明确的构造与评估接口。


\clearpage
\appendix
\section{补充设置、结果与候选审计}

\subsection{主实验运行设置}

表~\ref{tab:hyper} 根据 NELL23K 主运行配置和候选库元数据列出实际设置。Stage~2 从 Stage~1 adapter 初始化，优化候选级 InfoNCE 与生成损失。采样温度 $\tau_s$ 控制候选分布，候选温度 $\tau_c$ 控制对比目标。

\begin{table}[htbp]
\caption{NELL23K 主实验运行设置。数值对应训练种子 2026 的已完成运行。}
\label{tab:hyper}
\centering\small
\begin{tabular}{ll}
\toprule
字段 & 设置 \\
\midrule
基座模型 & Qwen2.5-7B-Instruct \\
候选数 $K$ & 9 \\
均匀混合系数 $\rho$ & $1/3$ \\
结构权重 $\alpha$ / 违例权重 $\beta$ & 1.0 / 1.0 \\
采样温度 $\tau_s$ & 0.25 \\
候选损失温度 $\tau_c$ & 1.0 \\
候选损失权重 $\lambda_{\mathrm{cand}}$ & 1.0 \\
LoRA rank / alpha & 4 / 8 \\
LoRA 目标模块 & \texttt{down\_proj}, \texttt{up\_proj}, \texttt{gate\_proj} \\
Stage~2 轮数 / 优化步数 & 10 / 240 \\
学习率 / 权重衰减 & $10^{-3}$ / $10^{-4}$ \\
单设备批大小 / 梯度累积步数 & 1 / 512 \\
计算精度 & bfloat16 \\
候选库随机种子 / 训练种子 & 2026 / 2026 \\
候选来源与过滤 & train-only \\
\bottomrule
\end{tabular}
\end{table}

\subsection{WN18RR 直接解码补充结果}
\label{sec:wn-exploratory}

WN18RR 的训练期候选库覆盖 2,709 条关系问题，每题包含九个互异候选，构造失败数为 0；来源检查记录为 train-only 且与任务问题对齐。在训练种子 2026 下，使用完整验证集选择候选损失权重 $\lambda_{\mathrm{cand}}=0.25$，获选模型在 3,134 条测试问题上的直接解码 EM 为 86.89\%，训练为 5 个 epoch、125 个优化步。

同一测试文件上的本地 B1 直接解码 EM 为 32.67\%。该本地基线与原始 GRIP 报告的五次运行均值 64.30\% 存在待核验的复现差异 \citep{feng2025grip}。WN18RR 的补充结果对应单个获选模型的直接解码评估，完整候选重排、同预算随机对照和多种子重复尚未完成。

\subsection{FB15K237-2 候选 InfoNCE 参考实验}

该参考实验展示独立配置下的候选 InfoNCE 表现。B1 与 InfoNCE-i0 均在同一本地测试文件的 10,000 条问题上评分；i0 使用单独的配对候选文件、LoRA rank 24、Stage~2 训练 5 个 epoch 和候选损失权重 0.1。表~\ref{tab:historical-fb} 单独标记该配置，本文的 SC-CL 主结果对应表~\ref{tab:nell-results}。

\begin{table}[htbp]
\caption{FB15K237-2 候选 InfoNCE 参考结果（\%，10,000 条测试问题）。InfoNCE-i0 为独立候选及训练配置，方法身份与 SC-CL 分别记录。}
\label{tab:historical-fb}
\centering\small
\begin{tabular}{lrrrr}
\toprule
方法 & 直接解码 EM & 重排 H@1 & MRR & H@3 \\
\midrule
\bone & 73.23 & 79.91 & 88.01 & 95.80 \\
InfoNCE-i0 & 88.69 & 91.62 & 95.41 & 99.20 \\
\bottomrule
\end{tabular}
\end{table}

\subsection{候选来源与真值状态审计}
\label{sec:audit-protocol}

候选库的来源检查核对训练划分、问题标识、目标关系、每题九个互异候选和训练已知关系过滤结果。文件哈希、脚本版本、候选评分及模型配置共同标识运行产物。来源合规与事实真值分别记录：训练已知关系冲突表示候选命中该实体对的训练关系，持出正例重叠表示候选命中持出划分中的已知正例，其余候选保留开放世界未判定状态。

实体替换与关系替换使用各自的候选空间。CoDEx 的人工核验实体替换负例用于其三元组判断设定 \citep{safavi2020codex}；固定实体对后的关系替换，按训练已知关系和持出正例分别核对真值状态。

\begin{table}[htbp]
\caption{候选来源与真值状态的审计规则。该表定义状态分类，实测来源检查结果见正文。}
\label{tab:audit-states}
\centering\small
\begin{tabular}{p{0.24\linewidth}p{0.65\linewidth}}
\toprule
状态 & 判定规则 \\
\midrule
训练已知关系冲突 & 候选属于当前实体对的训练已知关系集合 $T_{\mathrm{tr}}(h,t)$。 \\
持出正例重叠 & 候选对应的三元组出现在持出划分的已知正例中；该记录用于审计。 \\
开放世界未判定 & 候选未被已知正例或其他核验证据确认，真值维持未判定。 \\
候选不足 & 训练已知关系过滤后 $|P_q|<K$，记录该问题及原因。 \\
问题对齐异常 & 候选库与训练任务的问题标识、目标关系或候选文本不一致。 \\
\bottomrule
\end{tabular}
\end{table}

\subsection{候选库记录与构造伪代码}

每条候选库记录包含问题标识、数据划分、实体对、目标关系、候选关系和采样来源。逐候选记录包括结构相似度、违例度、utility、B1 候选分数及其与目标分数的差值；来源清单包含过滤划分、采样参数、随机种子、脚本版本和文件哈希。这些字段将候选采样结果关联到原始任务及其离线评分。

{\small
\begin{verbatim}
for q in train_questions:
    assert q.split == "train"
    P = train_relations - {q.gold} - known_train_relations(q.head, q.tail)
    if len(P) < K:
        record_insufficient_pool(q)
        continue
    for r in P:
        s = profile_similarity(phi[q.gold], phi[r])
        v = b1_score_violation(q, r)
        u[r] = alpha * s + beta * v
    candidates = sample_without_replacement(
        P, mixture_softmax(u, rho, tau_s), K=9)
    write_frozen_bank(q, candidates, manifest=provenance)
\end{verbatim}
}

候选库按训练来源限定、已知关系过滤和候选采样的顺序构建。来源检查核对问题标识与关系文本的一致性，固定候选库供整个 Stage~2 使用。

\subsection{数据、模型与产物记录}

主实验使用 NELL23K，补充结果分别使用 WN18RR 和 FB15K237-2；基座模型为 Qwen2.5-7B-Instruct \citep{qwen2024qwen25}，LoRA 用于参数高效训练 \citep{hu2022lora}。实验处理既有知识图谱关系问题，未新增个人敏感数据采集。预测、候选分数、训练配置、adapter 与来源清单保存在本地实验归档中，当前未提供公开下载地址。


\bibliographystyle{iclr2026_conference}
\bibliography{references}

\end{document}
